\section{Evaluation Plan}

% Track 3 REQUIREMENT (solicitation): "Describe how success will be measured. Include metrics, tools,
% and benchmarks for assessing progress and effectiveness. Ideally, the evaluation plan should
% incorporate testing and validation opportunities with existing users."
% Also a Track-3 review criterion: "a specific, actionable list of milestones and evaluation plan."

\noindent\textbf{Activity 1 (Discover).}
% TODO(pesose): metrics -- number/severity of confirmed vulnerabilities; code/edge coverage achieved
% on target libraries; reproduction rate; time-to-find vs. baselines. Benchmarks: known-bug corpora /
% injected-bug benchmarks for FreeRTOS libraries. \TODO{set concrete targets and name baselines}

\noindent\textbf{Activity 2 (Prevent/Harden).}
% TODO(pesose): metrics -- vulnerability classes eliminated by construction; hardening overhead
% (code size, RAM, run-time) on representative MCU targets (must stay within embedded budgets);
% SBOM coverage of the release; secure-dev practice conformance. \TODO{set concrete targets}

\noindent\textbf{Activity 3 (Detect/Recover).}
% TODO(pesose): metrics -- CI gate catch-rate on regression/injected bugs; time-from-disclosure-to-
% fix; fraction of fixes upstreamed and adopted; downstream-notification coverage via SBOM.

\noindent\textbf{Validation with existing users.}
A primary validation channel is the \ose-based drone/IoT testbed operated by collaborator Prof.\ Ranganathan (e.g., Crazyflie and ESP32-class platforms), on which discovered-and-hardened vulnerabilities are reproduced and the hardening's footprint/overhead is measured on real deployments.
% TODO(pesose): detail the validation protocol on Ranganathan's testbed; add any other FreeRTOS
% users/contributors who will validate (ties to the 3-5 required letters of collaboration).
\TODO{name additional validation partners / real deployments}
